\chapter{Penutup}
\label{chap:penutup}

\section{Kesimpulan}

Berdasarkan hasil analisis, perancangan, implementasi, dan evaluasi yang telah dilakukan, dapat disimpulkan beberapa hal sebagai berikut:

\begin{enumerate}
	\item Merancang dan mengimplementasikan kontrol akses granular pada DecMed dengan Macaroons sebagai token otorisasi. Batasan akses direpresentasikan melalui \textit{caveats}, sedangkan penegakan otoritas dilakukan pada \textit{Server} PRE sebelum data diakses.

	\item Segmentasi data RME diterapkan dengan memecah data \textit{off-chain} berdasarkan \textit{dataset category} dan \textit{function category}. Dengan struktur ini, data RME tidak lagi diperlakukan sebagai satu \textit{encrypted blob} besar, tetapi sebagai kumpulan segmen klinis. Metadata setiap segmen digunakan untuk menentukan segmen yang dapat diakses oleh aktor sesuai cakupan token.

	\item Atenuasi \textit{offline} diterapkan untuk mendukung pendelegasian akses antar tenaga kesehatan tanpa penerbitan token baru oleh pasien pada setiap perpindahan akses. Pemegang token membentuk token turunan secara \textit{offline} dengan menambahkan \textit{caveats} yang mempersempit cakupan token induk. Implementasi ini juga mencakup pencatatan delegasi, pembatasan kedalaman delegasi, validasi agar token turunan tidak memperluas hak akses token induk, serta revokasi berantai.
\end{enumerate}

\section{Saran}

Berdasarkan batasan dan hasil evaluasi Tugas Akhir ini, beberapa saran untuk pengembangan berikutnya adalah sebagai berikut:

\begin{enumerate}
	\item Cakupan skenario layanan klinis hingga pelayanan antar fasyankes, memungkinkan skenario nyata seperti rujukan antar fasyankes.

	\item Cakupan skenario layanan klinis hingga pelayanan gawat darurat. Mengimplementasikan data IGD dan mekanisme akses ketika pasien tidak sadar, memungkinkan skenario nyata gawat darurat.

	\item Integrasi dengan standar interoperabilitas data kesehatan dapat diteliti lebih lanjut. Segmentasi RME dapat dikembangkan untuk integrasi dengan SATUSEHAT atau pemetaan ke standar FHIR.

	\item Integrasi dan Pengujian dengan Sistem Rumah Sakit untuk mendukung penggunaan sistem dalam konteks nyata. Disesuaikan dengan tata laksana atau dokumen Pedoman Nasional Pelayanan Kesehatan.
\end{enumerate}